\setcounter{chaptercntr}{5}

\sectionbreak \section*{
	\gostTitleFont
	\redline
	\thechaptercntr \space
	ТЕХНИКО-ЭКОНОМИЧЕСКОЕ ОБОСНОВАНИЕ
}

\titlespace

\subsection*{ 
	\gostTitleFont
	\redline
	\thechaptercntr .\thesubchaptercntr \spc
	Исходные данные
} \addtocounter{subchaptercntr}{1}

\subtitlespace

{\gostFont
	
	
	\par \redline Задачей данного дипломного проекта является разработка серверной части приложения "Такси".
	\par \redline Разрабатываемое программное средство предназначено для обработки заказов, взаимодействия между пассажирами и водителями, хранения пользовательских данных, а также обеспечения работы клиентского приложения.
	\par \redline Разработка данного программного средства предусматривает проведение практически всех стадий проектирования (техническое задание, эскизный проект, технический проект, рабочий проект) и относится ко второй группе сложности.
	\par \redline Последовательность расчётов: 
	\par \redline – Расчёт объёма функций программного модуля;
	\par \redline – Расчёт полной себестоимости
	\par \redline – Расчёт цены и прибыли по программному продукту.
	
}

\subtitlespace

\subsection*{ 
	\gostTitleFont
	\redline
	\thechaptercntr .\thesubchaptercntr \spc
	Расчет объёма функций программного обеспечения
} \addtocounter{subchaptercntr}{1}

\subtitlespace

{\gostFont
	\par \redline Средой разработки серверной части приложения "Такси" является IntelliJ IDEA. Язык разработки: Java.
	\par \redline Общий объём ПО определяется по формуле (\thechaptercntr .\theformulacntr), исходя из объёма функций, реализуемых программой. 
	
	\formulaspace \par \redline 
	$V_0=\sum_{i=0}^{n}V_i$, 
	\hfill (\thechaptercntr .\theformulacntr) \redline
	\formulaspace \addtocounter{formulacntr}{1}
	
	\par \redline где $V_0$ {--} общий объём ПО;
	\par \redline \wherespace $V_i$ {--} объём функций ПО;
	\par \redline \wherespace n {--} общее число функций.
	\par \redline Расчёт общего объема ПР (количества строк исходного кода) предполагает определение объема по каждой функции. В том случае, когда на стадии технико-экономического обоснования проекта невозможно рассчитать точный объём функций, то данный объём может быть получен на основании прогнозируемой оценки имеющихся фактических данных по аналогичным проектам, выполненным ранее, или применением нормативов по ка-\\талогу функций. 
	
	\par \redline По каталогу функций на основании функций разрабатываемого ПО определяется общий объём ПО. Также на основе зависимостей от организационных и технологических условий, был скорректирован объём на основе экспертных оценок.
	
	\par \redline Уточнённый объём ПО определяется по формуле (\thechaptercntr .\theformulacntr).
	
	\formulaspace \par \redline 
	$V_y=\sum_{i=0}^{n}V_{yi}$,
	\hfill (\thechaptercntr .\theformulacntr) \redline
	\formulaspace \addtocounter{formulacntr}{1}
	
	\begin{tabular}{p{0,875cm}p{0,3cm}p{15,175cm}}
		& где & $V_{yi}$ {--} уточнённый объём отдельной функции в строках исходного кода. Перечень и объём функций ПО приведён в таблице \thechaptercntr .\thetablecntr .
	\end{tabular}
	
	\topTablespace
	{\begin{Center} 
			\par Таблица \thechaptercntr .\thetablecntr \spc {--} Перечень и объем функций программного обеспечения 
			
			\begin{tabular}{|c|p{6.5cm}|c|c|}
				\hline
				Номер   & Наименование (содержание)\centering                 & \multicolumn{2}{c|}{Объём функции строк исходного кода} \\ \cline{3-4}
				функции & `функции\centering& По каталогу $V_y$     & Уточнённый $V_{yi}$             \\ \hline
				102     & Контроль, предварительная обработка и ввод информации\centering & 490& 520\\ \hline
				202     & Формирование баз данных\centering                   & 1980& 2050\\ \hline
				203& Обработка наборов и записей базы данных\centering                   & 2370& 2280\\ \hline
				303     & Обработка файлов\centering                          & 1050& 900\\ \hline
				405     & Система настройки ПО\centering                      & 340& 200\\ \hline
				506     & Обработка ошибочных и сбойных ситуаций\centering    & 1540& 1320\\ \hline
				602     & Вспомогательные и сервисные программы\centering     & 470& 400\\ \hline
				703     & Расчет показателей\centering                        & 420& 350\\ \hline
				707     & Графический вывод результатов\centering             & 420& 300\\ \hline
				& ИТОГО\centering                                     & 9080& 8320\\ \hline
			\end{tabular}
			
			\botTablespace \end{Center}}
	\botTablespace
	
	\redline С учётом информации, указанной в таблице \thechaptercntr .\thetablecntr, уточнённый объём ПО составил 8320 строк кода вместо предполагаемого количества 9080 строк. \addtocounter{tablecntr}{1}
	
	\par
}

\subtitlespace

\subsection*{ 
	\gostTitleFont
	\redline
	\thechaptercntr .\thesubchaptercntr \spc
	Расчёт полной себестоимости программного модуля
} \addtocounter{subchaptercntr}{1}

\subtitlespace

{\gostFont
	
	\par \redline Стоимостная оценка программного средства у разработчика предполагает составление сметы затрат, которая включает следующие статьи расходов:
	
	\par \redline -	отчисления на социальные нужды ($\textrm{Р}_{\textrm{соц}}$);
	\par \redline -	материалы и комплектующие изделия ($\textrm{Р}_{\textrm{м}}$);
	\par \redline -	спецоборудование ($\textrm{Р}_{\textrm{с}}$);
	\par \redline -	машинное время ($\textrm{Р}_{\textrm{мв}}$);
	\par \redline -	расходы на научные командировки ($\textrm{Р}_{\textrm{нк}}$):
	\par \redline -	прочие прямые расходы ($\textrm{Р}_{\textrm{пр}}$);
	\par \redline -	накладные расходы ($\textrm{Р}_{\textrm{нр}}$);
	\par \redline -	затраты на освоение и сопровождение программного средства($\textrm{Р}_{\textrm{о}}$ и $\textrm{Р}_{\textrm{со}}$).
	
	\par \redline Полная себестоимость ($\textrm{С}_{\textrm{п}}$) разработки программного продукта рассчитывается как сумма расходов по всем статьям с учетом рыночной стоимости аналогичных продуктов.
	
	\par \redline Основной статьёй расходов на создание программного продукта является заработная плата проекта (основная и дополнительная) разработчиков (исполнителей) ($\textrm{ЗП}_{\textrm{осн}}$ + $\textrm{ЗП}_{\textrm{доп}}$), в число которых принято включать инженеров-программистов, руководителей проекта, системных архитекторов, дизайнеров, разработчиков баз данных, Web-мастеров и  других специалистов, необходимых для решения специальных задач в команде.
	
	\par \redline Расчёт заработной платы разработчиков программного продукта начинается с определения:
	\par \redline -	Продолжительности времени разработки ($\textrm{Ф}_{\textrm{рв}}$), которое устанавливается экспертным путем с учётом сложности, новизны программного продукта и фактически затраченного времени. В данном дипломном проекте $\textrm{Ф}_{\textrm{рв}}$ = 3 месяца.
	\par \redline -	Количества разработчиков программного обеспечения. В данном дипломном проекте три разработчика категории инженер-программист разряда 11. Заработная плата разработчиков определятся как сумма основной и дополнительной заработной платы всех исполнителей.
	
	\par \redline Заработная плата разработчиков определяется как сумма основной и \\дополнительной заработной платы всех исполнителей. 
	
	\par \redline Основная заработная плата $\textrm{ЗП}_{\textrm{осн}}$ каждого исполнителя определяется по формуле (\thechaptercntr .\theformulacntr):
	
	\formulaspace \par \redline 
	$\textrm{ЗП}_{\textrm{осн}} = \textrm{Т}_{\textrm{ст.1р}} \cdot \frac{\textrm{\FTwelwe Т}_{\textrm{\FNine к}}}{\textrm{\FTwelwe Ф}_{\textrm{\FNine эфф.р.в.}}} \cdot \textrm{Ф}_{\textrm{рв}} \cdot \textrm{К}_{\textrm{пр}}$ ,
	\hfill (\thechaptercntr .\theformulacntr) \redline
	\formulaspace 
	
	\begin{tabular}{p{0.875cm}p{0.3cm}p{15.175cm}}
		& где  & $\textrm{Т}_{\textrm{ст.1р}}$ {--} размер базовой ставки, на текущий момент актуальное значение которой 297 бел.руб.; \\
		& 	   & $\textrm{Т}_{\textrm{к}}$ {--} тарифный коэффициент согласно разряду исполнителя; \\
		&	   & $\textrm{Ф}_{\textrm{эфф.р.в.}}$ {--} среднее количество рабочих дней; \\
		&	   & $\textrm{Ф}_{\textrm{рв}}$ {--} фонд рабочего времени исполнителя (продолжительность разработки программного модуля), дни; \\
		&	   & $\textrm{К}_{\textrm{пр}}$ {--} коэффициент премии, $\textrm{К}_{\textrm{пр}}$ = 1,5.
	\end{tabular}
	
	\par \redline Тарифный коэффициент согласно 11 разряда инженера-программиста $\textrm{Т}_{\textrm{к}}$ = 1,91. Продолжительность разработки программного продукта – 70 дней, количество рабочих дней в месяце на текущий год составляет 22 дня. Рассчитаем основную заработную плату инженера-программиста, согласно формуле (\thechaptercntr .\theformulacntr): \addtocounter{formulacntr}{1}
	
	\formulaspace \par \redline 
	$\textrm{ЗП}_{\textrm{осн}} \textrm{ = 297} \cdot \frac{\textrm{\FTwelwe 1,91}}{\textrm{\FTwelwe 22}} \cdot \textrm{70} \cdot \textrm{1,5 = 2707,43 (бел.руб.)}$ 
	\formulaspace 
	
	\par \redline Дополнительная заработная плата $\textrm{ЗП}_{\textrm{доп}}$ каждого исполнителя рассчитывается от основной заработной платы по формуле (\thechaptercntr .\theformulacntr).
	
	\formulaspace \par \redline 
	$\textrm{ЗП}_{\textrm{доп}} = \textrm{ЗП}_{\textrm{осн}} \cdot \frac{\textrm{\FTwelwe Н}_{\textrm{\FNine доп.зп}}}{\textrm{\FTwelwe 100\%}}$,
	\hfill (\thechaptercntr .\theformulacntr) \redline
	\formulaspace
	
	\par \redline где $\textrm{Н}_{\textrm{доп.зп}}$ {--} надбавка дополнительной заработной платы, равная 15\%.
	
	\par \redline Рассчитаем дополнительную заработную плату инженера-программиста, согласно формуле (5.4):
	
	\par \redline Рассчитаем заработную плату исполнителей проекта и результаты занесем в таблицу (\thechaptercntr .\thetablecntr): \addtocounter{formulacntr}{1}
	
	\formulaspace \par \redline 
	$\textrm{ЗП}_{\textrm{доп}} = \textrm{2707,43} \cdot \frac{\textrm{\FTwelwe 15\%}}{\textrm{\FTwelwe 100\%}} \textrm{ = 406,11 (бел.руб.)}$
	\formulaspace
	
	\par \redline Результаты вычислений внесём в таблицу \thechaptercntr .\thetablecntr.
	
	\topTablespace
	{\begin{Center} 
			\par Таблица \thechaptercntr .\thetablecntr \spc {--} Расчет заработной платы
			
			\begin{tabular}{|c|c|c|c|c|c|c|c|}
				\hline
				Категория  & \multirow{2}{*}{Разряд } & \multirow{2}{*}{$\textrm{T}_{\textrm{к}}$} & \multirow{2}{*}{$\textrm{Ф}_{\textrm{рв}}$} & \multirow{2}{*}{$\textrm{К}_{\textrm{пр}}$} & \multicolumn{3}{c|}{Заработная плата, бел.руб.} \\ \cline{6-8}
				работников & & & & & осн. & доп. & всего \\ \hline
				Инженер-   & \multirow{2}{*}{11} & \multirow{2}{*}{1,91} & \multirow{2}{*}{70} & \multirow{2}{*}{1,5}& \multirow{2}{*}{2707,43} & \multirow{2}{*}{406,11} & \multirow{2}{*}{3113,54} \\
				программист& & & & & & & \\ \hline
				Инженер-   & \multirow{2}{*}{11} & \multirow{2}{*}{1,91} & \multirow{2}{*}{70} & \multirow{2}{*}{1,5} & \multirow{2}{*}{2707,43} & \multirow{2}{*}{406,11} & \multirow{2}{*}{3113,54} \\
				программист& & & & & & & \\ \hline
				Инженер-   & \multirow{2}{*}{11} & \multirow{2}{*}{1,91} & \multirow{2}{*}{70} & \multirow{2}{*}{1,5} & \multirow{2}{*}{2707,43} & \multirow{2}{*}{406,11} & \multirow{2}{*}{3113,54} \\
				программист& & & & & & & \\ \hline
				Итого   & {--} & {--} & {--} & {--} & 8122,29 & 1218,33 & 9340,62 \\ \hline
	\end{tabular} \end{Center}} \addtocounter{tablecntr}{1}
	\botTablespace
	
	\par \redline Таким образом, как видно из таблицы, заработная плата инженеров-программистов составляет 9340,62 бел. руб.
	
	\par \redline  Отчисления на социальные нужды $\textrm{Р}_{\textrm{соц}}$ определяются по формуле (\thechaptercntr .\theformulacntr) в соответствии с действующим законодательством по нормативу (29\% – отчисления в ФСЗН + 6\% – отчисления по обязательному страхованию).
	
	\formulaspace \par \redline 
	$\textrm{Р}_{\textrm{соц}} = (\textrm{ЗП}_{\textrm{осн}} + \textrm{ЗП}_{\textrm{доп}}) \cdot \frac{\textrm{\FTwelwe 35\%}}{\textrm{\FTwelwe 100\%}}$
	\hfill (\thechaptercntr .\theformulacntr) \redline
	\formulaspace \addtocounter{formulacntr}{1}
	
	\par \redline Рассчитанные отчисления на социальные нужды:
	
	\formulaspace \par \redline 
	$\textrm{Р}_{\textrm{соц}} = \textrm{(8122,29 + 1218,33)} \cdot \frac{\textrm{\FTwelwe 35\%}}{\textrm{\FTwelwe 100\%}} \textrm{ = 3269,22 (бел.руб.).}$
	\formulaspace 
	\par \redline Расходы по статье cпецоборудование ($\textrm{Р}_{\textrm{c}}$) включает затраты на приобретение технических и программных средств специального назначения, необходимых для разработки методического пособия, включая расходы на проектирование, изготовление, отладку и другое. 
	\par \redline В данном дипломном проекте для разработки приобретение какого-либо спецоборудования не предусматривалось. Так как спецоборудование не было приобретено, данная статья не рассчитывается. 
	\par \redline Расходы на материалы и комплектующие $\textrm{Р}_{\textrm{м}}$ отражают расходы на магнитные носители, бумагу, красящие ленты и другие материалы, необходимые для разработки программного продукта. Норма расхода материалов в суммарном выражении определяется в процентах к основной заработной плате (\thechaptercntr .\theformulacntr).
	
	\formulaspace \par \redline 
	$\textrm{Р}_{\textrm{м}} = \textrm{ЗП}_{\textrm{осн}} \cdot \frac{\textrm{\FTwelwe Н}_{\textrm{\FNine мз}}}{\textrm{\FTwelwe 100\%}}$,
	\hfill (\thechaptercntr .\theformulacntr) \redline
	\formulaspace \addtocounter{formulacntr}{1}
	
	\par \redline где $\textrm{Н}_{\textrm{мз}}$ {--} норма расхода материалов от основной заработной платы, в процентах.
	
	\par \redline Рассчитаем расходы на материалы и комплектующий, приняв $\textrm{Н}_{\textrm{мз}}$ равным 4\%:
	
	\formulaspace \par \redline 
	$\textrm{Р}_{\textrm{м}} = \textrm{8122,29} \cdot \frac{\textrm{\FTwelwe 4\%}}{\textrm{\FTwelwe 100\%}} \textrm{ = 324,89 (бел.руб.)}$ 
	\formulaspace 
	
	\par \redline Расходы на машинное время $\textrm{Р}_{\textrm{мв}}$ включают оплату машинного времени, необходимого для разработки и отладки программного продукта. Они определяются в машино-часах по нормативам на 100 строк исходного кода машинного времени. $\textrm{Р}_{\textrm{мв}}$ определяется по формуле (\thechaptercntr .\theformulacntr).
	
	\formulaspace \par \redline 
	$\textrm{Р}_{\textrm{мв}} = \textrm{Ц}_{\textrm{мвi}} \cdot \frac{\textrm{\FTwelwe V}_{\textrm{\FNine 0}}}{\textrm{\FTwelwe 100}} \cdot \textrm{Н}_{\textrm{мв}}$,
	\hfill (\thechaptercntr .\theformulacntr) \redline
	\formulaspace \addtocounter{formulacntr}{1}
	
	\begin{tabular}{p{0,875cm}p{0,3cm}p{15,175cm}}
		& где  & $\textrm{Ц}_{\textrm{мвi}}$ {--} цена одного машинного часа (7 бел. руб.); \\
		& 	   & $\textrm{V}_{\textrm{0}}$ {--} уточнённый общий объём машинного кода; \\
		&	   & $\textrm{Н}_{\textrm{мв}}$ {--} норматив расхода машинного времени на отладку 100 строк кода в маши- но-часах. Принимается в размере 0,7. 
	\end{tabular}
	
	\par \redline Рассчитаем расходы на машинное время:
	
	\formulaspace \par \redline 
	$\textrm{Р}_{\textrm{мв}} = \textrm{7} \cdot \frac{\textrm{\FTwelwe 8320}}{\textrm{\FTwelwe 100}} \cdot \textrm{0,7 = 407,68 (бел. руб.)}$
	\formulaspace
	
	\par \redline Расходы по статье "Научные командировки" $\textrm{Р}_{\textrm{нк}}$ берутся либо по смете научных командировок, разрабатываемой на предприятии, либо в процентах от основной заработной платы исполнителей (10-15\%). Так как в данном проекте научные командировки не предусмотрены, данная статья не рассчитывается. 
	
	\par \redline Прочие затраты $\textrm{Р}_{\textrm{пр}}$ включают затраты на приобретение специальной научно-тех- нической информации и специальной литературы. Определяются по нормативу в процентах к основной заработной плате исполнителей. Так как специальная научно-техническая информация и специальная литература не приобреталась, то данная статья не рассчитывается. 
	
	\par \redline Затраты на накладные расходы $\textrm{Р}_{\textrm{нр}}$ связаны с содержанием вспомогательных хозяйств, опытных производств, а также с расходами на общехозяйственные нужды. Определяется по нормативу в процентах к основной заработной плате по формуле (\thechaptercntr .\theformulacntr).
	
	\formulaspace \par \redline 
	$\textrm{Р}_{\textrm{нр}} = \frac{\textrm{\FTwelwe Н}_{\textrm{\FNine нр}}}{\textrm{\FTwelwe 100\%}} \cdot \textrm{ЗП}_{\textrm{осн}}$,
	\hfill (\thechaptercntr .\theformulacntr) \redline
	\formulaspace 
	
	\par \redline где $\textrm{Н}_{\textrm{нр}}$ {--} норматив накладных расходов.
	
	\par \redline В данном дипломном проекте норматив накладных расходов равен 75\%, поэтому затраты на накладные расходы согласно формуле (\thechaptercntr .\theformulacntr) равны: \addtocounter{formulacntr}{1}
	
	\formulaspace \par \redline 
	$\textrm{Р}_{\textrm{нр}} = \frac{\textrm{\FTwelwe 75\%}}{\textrm{\FTwelwe 100\%}} \cdot \textrm{8122,29 =  6091,72(бел.руб.).}$
	\formulaspace
	
	\par \redline Сумма вышеперечисленных расходов по статьям на программный продукт служит исходной базой для расчёта затрат на освоение и сопровождение программного продукта. Они рассчитываются по формуле (\thechaptercntr .\theformulacntr):
	
	\formulaspace \par \redline 
	$\textrm{СЗ = } \textrm{ЗП}_{\textrm{осн}} + \textrm{ЗП}_{\textrm{доп}} + \textrm{Р}_{\textrm{соц}} + \textrm{Р}_{\textrm{м}} + \textrm{Р}_{\textrm{с}} + \textrm{Р}_{\textrm{мв}} + \textrm{Р}_{\textrm{нк}} + \textrm{Р}_{\textrm{пр}} + \textrm{Р}_{\textrm{нр}}$
	\hfill (\thechaptercntr .\theformulacntr) \redline
	\formulaspace \addtocounter{formulacntr}{1}
	
	\par \redline тогда
	
	\formulaspace \par \redline 
	$\textrm{СЗ = 8122,29 + 1218,33 + 3269,22 + 324,89 + 407,68 + 6091,72 = 19434,13 (бел.руб.).}$
	\formulaspace 
	
	\par \redline Организация-разработчик участвует в освоении программного продукта и несёт соответствующие затраты, на которые составляется смета, оплачиваемая заказчиком по договору. Затраты на освоение $\textrm{Р}_{\textrm{о}}$ определяются по установленному нормативу от суммы затрат по формуле (\thechaptercntr .\theformulacntr).
	
	\formulaspace \par \redline 
	$\textrm{Р}_{\textrm{о}} = \textrm{СЗ} \cdot \frac{\textrm{\FTwelwe Н}_{\textrm{\FNine о}}}{\textrm{\FTwelwe 100\%}}$,
	\hfill (\thechaptercntr .\theformulacntr) \redline
	\formulaspace \addtocounter{formulacntr}{1}
	
	\begin{tabular}{p{0,875cm}p{0,3cm}p{15,175cm}}
		& где  & $\textrm{СЗ}$ {--} сумма вышеперечисленных расходов по статьям на разработку программного продукта; \\
		& 	   & $\textrm{Н}_{\textrm{о}}$ {--} установленный норматив затрат на освоение. Для данного дипломного проекта принимается равным 8\%. 
	\end{tabular}
	
	\par \redline Рассчитаем затраты на освоение продукта:
	
	\formulaspace \par \redline 
	$\textrm{Р}_{\textrm{о}} = \textrm{19434,13} \cdot \frac{\textrm{\FTwelwe 8\%}}{\textrm{\FTwelwe 100\%}} \textrm{ = 1554,73 (бел.руб.)}$
	\formulaspace 
	
	\par \redline Организация-разработчик осуществляет сопровождение программного продукта и несёт расходы, которые оплачиваются заказчиком в соответствии с договором и сметой на сопровождение. Расходы на сопровождение $\textrm{Р}_{\textrm{со}}$ рассчитываются по формуле (\thechaptercntr .\theformulacntr).
	
	\formulaspace \par \redline 
	$\textrm{Р}_{\textrm{со}} = \textrm{СЗ} \cdot \frac{\textrm{\FTwelwe Н}_{\textrm{\FNine со}}}{\textrm{\FTwelwe 100\%}}$,
	\hfill (\thechaptercntr .\theformulacntr) \redline
	\formulaspace \addtocounter{formulacntr}{1}
	
	\begin{tabular}{p{0,875cm}p{0,3cm}p{15,175cm}}
		& где  & $\textrm{Н}_{\textrm{со}}$ {--} установленный норматив затрат на сопровождение программного продукта. Для данного дипломного проекта принимается равным 7\%.
	\end{tabular}
	
	\par \redline Рассчитаем расходы на сопровождение:
	
	\formulaspace \par \redline 
	$\textrm{Р}_{\textrm{со}} = \textrm{19434,13} \cdot \frac{\textrm{\FTwelwe 7\%}}{\textrm{\FTwelwe 100\%}} \textrm{ = 1360,39 (бел.руб.)}$
	\formulaspace 
	
	\par \redline Полная себестоимость (СП) разработки программного продукта рассчитывается как сумма расходов по всем статьям. Она определяется по формуле (\thechaptercntr .\theformulacntr).
	
	\formulaspace \par \redline 
	$\textrm{СП} = \textrm{СЗ} + \textrm{Р}_{\textrm{о}} + \textrm{Р}_{\textrm{со}}$
	\hfill (\thechaptercntr .\theformulacntr) \redline
	\formulaspace \addtocounter{formulacntr}{1}
	
	\par \redline Рассчитанная полная себестоимость продукта:
	
	\formulaspace \par \redline 
	$\textrm{СП = 19434,13 + 1554,73 + 1360,39 = 22349,25(бел.руб.)}$
	\formulaspace 
	
	\par \redline Результаты вычислений занесём в таблицу \thechaptercntr .\thetablecntr.
	
	\topTablespace
	{\begin{Center}
			\par Таблица \thechaptercntr .\thetablecntr \spc {--} Себестоимость программного продукта
			
			\begin{tabular}{|c|c|c|}
				\hline
				Наименование статей затрат & Норматив, \% & Сумма затрат, бел.руб. \\ \hline
				Заработная плата, всего & {--} & 9340,62\\ \hline
				Основная заработная плата & {--} & 8122,29\\ \hline
				Дополнительная заработная плата & {--} & 1218,33\\ \hline
				Отчисления на социальные нужды & 35 & 3269,22\\ \hline
				Спецоборудование & Не применялось & {--} \\ \hline
				Материалы и комплектующие изделия & 4 & 324,89\\ \hline
				Машинное время & {--} & 407,68\\ \hline
				Научные командировки & Не планировались & {--} \\ \hline
				Прочие затраты & Не применялись & {--} \\ \hline
				Накладные расходы & 75 & 6091,72\\ \hline
				Сумма затрат & {--} & 19434,13\\ \hline
				Затраты на освоение & 8 & 1554,73\\ \hline
				Затраты на сопровождение & 7 & 1360,39\\ \hline
				Полная себестоимость & {--} & 22349,25\\ \hline
	\end{tabular} \end{Center}} \addtocounter{tablecntr}{1} 
	\botTablespace
	
	\par \redline В результате всех расчётов полная себестоимость программного продукта составила 22349,25 бел.руб. 
}


\subtitlespace

\subsection*{ 
	\gostTitleFont
	\redline
	\thechaptercntr .\thesubchaptercntr \spc
	Расчет цены и прибыли по программному продукту
} \addtocounter{subchaptercntr}{1}

\subtitlespace

{\gostFont
	
	\par \redline Для определения цены программного продукта необходимо рассчитать плановую прибыль П, которая рассчитывается по формуле (\thechaptercntr .\theformulacntr). 
	
	\formulaspace \par \redline 
	$\textrm{П} = \textrm{СП} \cdot \frac{\textrm{\FTwelwe R}}{\textrm{\FTwelwe 100\%}}$,
	\hfill (\thechaptercntr .\theformulacntr) \redline
	\formulaspace \addtocounter{formulacntr}{1}
	
	\par \redline где $\textrm{СП}$ {--} полная себестоимость программного модуля, бел. руб; 
	\par \redline \wherespace $\textrm{R}$ {--} уровень рентабельности программного модуля. 
	
	\par \redline Рассчитаем прибыль от реализации:
	
	\formulaspace \par \redline 
	$\textrm{П = 22349,25} \cdot \frac{\textrm{\FTwelwe 23\%}}{\textrm{\FTwelwe 100\%}} \textrm{ = 5140,33(бел.руб)}$
	\formulaspace
	
	\par \redline После расчета прибыли от реализации по формуле (\thechaptercntr .\theformulacntr) определяется прогнозируемая цена программного продукта без налогов $\textrm{Ц}_{\textrm{п}}$.
	
	\formulaspace \par \redline 
	$\textrm{Ц}_{\textrm{п}} = \textrm{СП} + \textrm{П}$
	\hfill (\thechaptercntr .\theformulacntr) \redline
	\formulaspace \addtocounter{formulacntr}{1}
	
	\par \redline Рассчитаем цену программного продукта без налогов:
	
	\formulaspace \par \redline 
	$\textrm{Ц}_{\textrm{п}} = \textrm{22349,25 + 5140,33 = 27489,58 (бел.руб.)}$
	\formulaspace
	
	\par \redline Отпускная цена $\textrm{Ц}_{\textrm{о}}$ (цена реализации) программного продукта включает налог на добавленную стоимость и рассчитывается по формуле (\thechaptercntr .\theformulacntr).
	
	\formulaspace \par \redline 
	$\textrm{Ц}_{\textrm{о}} = \textrm{Ц}_{\textrm{п}} + \textrm{НДС}_{\textrm{пп}}$,
	\hfill (\thechaptercntr .\theformulacntr) \redline
	\formulaspace \addtocounter{formulacntr}{1}
	
	\par \redline где $\textrm{НДС}_{\textrm{пп}}$ {--} налог на добавленную стоимость для программного продукта.
	
	\par \redline Для данного программного продукта $\textrm{НДС}_{\textrm{пп}}$ рассчитывается по формуле (\thechaptercntr .\theformulacntr).
	
	\formulaspace \par \redline 
	$\textrm{НДС}_{\textrm{пп}} = \textrm{Ц}_{\textrm{п}} \cdotp \frac{\textrm{\FTwelwe НДС}}{\textrm{\FTwelwe 100\%}}$,
	\hfill (\thechaptercntr .\theformulacntr) \redline
	\formulaspace \addtocounter{formulacntr}{1}
	
	\begin{tabular}{p{0,875cm}p{0,3cm}p{15,175cm}}
		& где & $\textrm{НДС}$ {--} налог на добавленную стоимость. В настоящее время он составляет 20\%.
	\end{tabular}
	
	\par \redline Рассчитаем $\textrm{НДС}_{\textrm{пп}}$:
	
	\formulaspace \par \redline 
	$\textrm{НДС}_{\textrm{пп}} = \textrm{27489,58} \cdot \frac{\textrm{\FTwelwe 20\%}}{\textrm{\FTwelwe 100\%}} \textrm{ = 5497,92 (бел.руб.)}$
	\formulaspace
	
	\par \redline Рассчитаем $\textrm{Ц}_{\textrm{о}}$:
	
	\formulaspace \par \redline 
	$\textrm{Ц}_{\textrm{о}} = \textrm{27489,58 + 5497,92 = 32987,50 (бел.руб.)}$
	\formulaspace
	
	\par \redline Прибыль от реализации программного продукта за вычетом налога на прибыль является чистой прибылью. Чистая прибыль остаётся организации-разработчику и представляет собой экономический эффект от создания нового программного продукта. Она рассчитывается по формуле (\thechaptercntr .\theformulacntr).
	
	\formulaspace \par \redline 
	$\textrm{ПЧ} = \textrm{П} \cdot (\textrm{1} - \frac{\textrm{\FTwelwe Н}_{\textrm{\FNine п}}}{\textrm{\FTwelwe 100\%}})$,
	\hfill (\thechaptercntr .\theformulacntr) \redline
	\formulaspace \addtocounter{formulacntr}{1}
	
	\par \redline где $\textrm{Н}_{\textrm{п}}$ {--} ставка налога на прибыль. В настоящее время он равен 20\%. 
	
	\par \redline Рассчитаем чистую прибыль:
	
	\formulaspace \par \redline 
	$\textrm{ПЧ} = \textrm{5140,33} \cdot (\textrm{1} - \frac{\textrm{\FTwelwe 20\%}}{\textrm{\FTwelwe 100\%}}) \textrm{ = 4112,26 {бел.руб.}}$
	\formulaspace
	
	\par \redline Результаты расчётов цены и прибыли по программному продукту сведены в таблицу \thechaptercntr .\thetablecntr.
	
	\topTablespace
	{\begin{Center}
			\par Таблица \thechaptercntr .\thetablecntr \spc {--} Расчёт отпускной цены и чистой прибыли программного модуля
			
			\begin{tabular}{|c|c|c|}
				\hline
				Наименование статей затрат & Норматив, \% & Сумма затрат, бел.руб. \\ \hline
				Полная себестоимость & {--} & 22349,25\\ \hline
				Прибыль & 23 & 5140,33\\ \hline
				Цена без НДС & {--} & 27489,58\\ \hline
				НДС & 20 & 5497,92\\ \hline
				Отпускная цена & {--} & 32987,50\\ \hline
				Налог на прибыль & 20 & 1028,07\\ \hline
				Чистая прибыль & {--} & 4112,26\\ \hline
	\end{tabular} \end{Center}} \addtocounter{tablecntr}{1}
	\botTablespace
	
	\par \redline В ходе произведенных расчетов определены основные экономические показатели: полная себестоимость программного продукта составила 22349,25 бел.руб.; прогнозируемая цена без НДС – 27489,58 бел.руб.; отпускная цена программного продукта – 32987,50 бел.руб.; чистая прибыль – 4112,26 бел.руб.
	
	\par \redline Разработанный программный продукт имеет малое количество конкурентов с более высокими ценами на их услуги. Таким образом, рассчитанная отпускная цена на программный продукт, разрабатываемой в рамках данного дипломного проекта, является конкурентоспособной. При расчете цены учтены отчисления в фонд социальной защиты, а также налоги, необходимые к уплате. К конечному итогу получаем окончательную цену продукта, равную 32987,50 белорусских рублей.
	
}

\setcounter{subchaptercntr}{1}
\setcounter{formulacntr}{1}
\setcounter{imagecntr}{1}